\documentclass[12pt]{report}

% ─── Paquetes básicos ────────────────────────────────────────────────────────
\usepackage[a4paper, top=2.5cm, bottom=2.5cm, left=3cm, right=2.5cm]{geometry}
\usepackage[spanish]{babel}
\usepackage[utf8]{inputenc}
\usepackage[T1]{fontenc}
\usepackage{graphicx}
\usepackage{float}
\usepackage{hyperref}
\usepackage{listings}
\usepackage{xcolor}
\usepackage{booktabs}
\usepackage{array}
\usepackage{parskip}
\usepackage{enumitem}
\usepackage{caption}
\usepackage{subcaption}

% ─── Configuración de hiperlinks ─────────────────────────────────────────────
\hypersetup{
    colorlinks=true,
    linkcolor=black,
    urlcolor=blue,
    citecolor=black
}

% ─── Configuración de listings (código fuente) ───────────────────────────────
\definecolor{codegray}{rgb}{0.95,0.95,0.95}
\definecolor{codeblue}{rgb}{0.1,0.2,0.6}
\definecolor{codecomment}{rgb}{0.3,0.5,0.3}

\lstdefinestyle{pythonstyle}{
    backgroundcolor=\color{codegray},
    commentstyle=\color{codecomment},
    keywordstyle=\color{codeblue},
    numberstyle=\tiny\color{gray},
    stringstyle=\color{orange!70!black},
    basicstyle=\ttfamily\footnotesize,
    breakatwhitespace=false,
    breaklines=true,
    captionpos=b,
    keepspaces=true,
    numbers=left,
    numbersep=5pt,
    showspaces=false,
    showstringspaces=false,
    tabsize=4,
    frame=single
}
\lstset{style=pythonstyle}

% ─── Metadatos del documento ─────────────────────────────────────────────────
\title{
    \textbf{Reporte de Avance Técnico} \\[0.5em]
    \large Trabajo Terminal 2 \\[0.3em]
    \normalsize Sistema de Visión Computacional para el Análisis \\
    Automatizado del Comportamiento de Roedores en el EPM \\[1em]
    \small Instituto Politécnico Nacional \\
    Escuela Superior de Cómputo (ESCOM) \\
    Centro de Apoyo a la Titulación (CATT)
}
\author{Trabajo Terminal TT\_Ratones\_2026}
\date{Abril 2026}

% ─────────────────────────────────────────────────────────────────────────────

\begin{document}

\maketitle
\tableofcontents
\clearpage

% =============================================================================
\chapter{Descripción General del Sistema}
% =============================================================================

\section{Contexto y motivación}

El análisis del comportamiento animal en entornos estandarizados constituye
una herramienta fundamental en la investigación neurocientífica y farmacológica.
El Elevated Plus Maze (EPM) es uno de los paradigmas más utilizados para evaluar
respuestas de ansiedad en roedores: el animal es colocado en un laberinto en
forma de cruz elevado sobre el suelo, con dos brazos abiertos y dos brazos
cerrados. La proporción de tiempo que el animal pasa en cada zona, su velocidad
de desplazamiento y los patrones motores que exhibe proporcionan métricas
objetivas del estado conductual del individuo.

Históricamente, este análisis se ha realizado de forma manual o mediante
software comercial de costo elevado. Ambas alternativas presentan limitaciones
importantes: la primera introduce sesgo observacional y consume un tiempo
considerable de trabajo humano especializado; la segunda restringe el acceso a
grupos de investigación con recursos limitados y dificulta la reproducibilidad
de los análisis.

\section{Propósito del sistema}

El sistema desarrollado en este trabajo terminal tiene como objetivo automatizar
el proceso completo de análisis conductual para experimentos EPM, desde la
recepción del video original hasta la generación de métricas cuantitativas
validadas. El sistema está diseñado para reducir la intervención manual al
mínimo indispensable, mantener trazabilidad completa de cada experimento y
producir resultados reproducibles entre sesiones.

En su estado actual, el sistema implementa el flujo completo de trabajo:

\begin{enumerate}[leftmargin=*, label=\arabic*.]
    \item Ingesta y validación del video de experimento.
    \item Extracción de keypoints corporales del roedor mediante DeepLabCut SuperAnimal.
    \item Filtrado espacial de poses mediante un detector YOLO entrenado localmente.
    \item Cómputo de características conductuales con SimBA.
    \item Clasificación de conductas: thigmotaxis y grooming.
    \item Generación de un video multimodal con anotaciones superpuestas.
    \item Persistencia de resultados en base de datos relacional.
    \item Visualización de métricas, trayectorias y estadísticas en un dashboard interactivo.
\end{enumerate}

\section{Alcance actual}

El sistema procesa videos de entre 5 y 10 minutos de duración en formato
estándar (MP4, AVI), capturados desde una cámara cenithal fija sobre el aparato
EPM. El escenario experimental está controlado en cuanto a iluminación y escala.
La interfaz de usuario opera en la web a través de una aplicación multipágina
desarrollada con Streamlit, accesible desde la red local del laboratorio.

La fase actual del proyecto se encuentra en una etapa de consolidación del
sistema operativo y en la apertura de una segunda línea de investigación
orientada a sustituir DeepLabCut como backend de extracción de pose por un
modelo YOLO Pose de última generación, con el objetivo principal de reducir
drásticamente los tiempos de inferencia.

% =============================================================================
\chapter{Arquitectura del Sistema}
% =============================================================================

\section{Visión general}

El sistema se organiza en tres capas funcionales claramente diferenciadas:
la capa de aplicación, la capa de lógica y soporte, y la capa de ejecución
del pipeline. Esta separación permite que la interfaz gráfica opere de forma
independiente de los procesos computacionales intensivos, los cuales se ejecutan
en segundo plano con retroalimentación en tiempo real al usuario.

\begin{figure}[H]
    \centering
    % Si se agrega un diagrama, descomentar y reemplazar el nombre del archivo:
    % \includegraphics[width=0.85\textwidth]{arquitectura_sistema.png}
    \fbox{\parbox{0.85\textwidth}{
        \centering\vspace{1em}
        \textit{[Diagrama de arquitectura -- insertar imagen]}
        \vspace{1em}
    }}
    \caption{Arquitectura de tres capas del sistema TT Ratones 2026.}
    \label{fig:arquitectura}
\end{figure}

\section{Capa de aplicación}

La interfaz de usuario está construida con Streamlit y se estructura como una
aplicación multipágina. Cada página corresponde a una etapa lógica del flujo
de trabajo experimental:

\begin{table}[H]
    \centering
    \begin{tabular}{@{}lll@{}}
        \toprule
        \textbf{Archivo} & \textbf{Página} & \textbf{Función principal} \\
        \midrule
        \texttt{Home.py}                     & Inicio            & Presentación y acceso al sistema \\
        \texttt{pages/00\_Login.py}           & Autenticación     & Inicio de sesión y verificación OTP \\
        \texttt{pages/01\_Ingesta\_de\_Video.py} & Ingesta        & Carga y validación del video \\
        \texttt{pages/02\_Keypoints.py}       & Keypoints         & Extracción de pose (DLC / YOLO Pose) \\
        \texttt{pages/03\_Configuracion\_Zonas.py} & Zonas       & Definición de ROIs en el EPM \\
        \texttt{pages/04\_Analisis\_Final.py}  & Análisis          & Ejecución del pipeline completo \\
        \texttt{pages/05\_Resultados\_y\_Estadisticas.py} & Resultados & Dashboard de métricas \\
        \texttt{pages/98\_Perfil.py}           & Perfil            & Gestión de cuenta de usuario \\
        \texttt{pages/99\_Admin\_Panel.py}     & Administración    & Gestión de usuarios y sistema \\
        \bottomrule
    \end{tabular}
    \caption{Páginas de la aplicación Streamlit y su función.}
    \label{tab:paginas}
\end{table}

\section{Capa de lógica y soporte}

El núcleo técnico reside en el directorio \texttt{src/} e incluye:

\begin{itemize}[leftmargin=*]
    \item \texttt{src/config.py}: configuración centralizada de rutas y parámetros del proyecto.
    \item \texttt{src/session\_utils.py}: gestión del estado de sesión entre páginas de Streamlit.
    \item \texttt{src/ui\_components.py}: componentes de interfaz reutilizables (splash, banners, etc.).
    \item \texttt{src/auth.py}: autenticación de usuarios con soporte de OTP por correo electrónico.
    \item \texttt{src/email\_utils.py}: envío de notificaciones y tokens vía SMTP.
    \item \texttt{src/analysis\_logic.py}: lógica base para interpretación de resultados conductuales.
    \item \texttt{src/simba\_roi\_bridge.py}: sincronización de zonas del EPM con el formato de regiones
          de interés que espera SimBA.
    \item \texttt{src/db/}: capa de acceso a base de datos (SQLite), incluyendo gestión de experimentos
          e historial.
    \item \texttt{src/security\_logger.py}: registro de eventos de seguridad y auditoría.
\end{itemize}

\section{Capa de ejecución del pipeline}

Los procesos computacionales intensivos se implementan como scripts independientes
invocados por el orquestador principal:

\begin{itemize}[leftmargin=*]
    \item \texttt{src/scripts/run\_behavior\_pipeline.py}: orquestador central del pipeline.
    \item \texttt{src/scripts/run\_superanimal.py}: extracción de pose con DLC SuperAnimal.
    \item \texttt{src/scripts/apply\_dlc\_bbox\_constraint.py}: filtrado de keypoints fuera del bounding box YOLO.
    \item \texttt{src/scripts/compute\_simba\_features.py}: cómputo de características para SimBA.
    \item \texttt{src/scripts/simba\_inference.py}: inferencia conductual con modelos SimBA entrenados.
    \item \texttt{src/scripts/generar\_video\_prediccion.py}: generación del video multimodal de resultados.
    \item \texttt{src/scripts/run\_with\_live\_log.py}: envoltorio para ejecución con transmisión de logs en tiempo real.
    \item \texttt{src/scripts/render\_dlc\_keypoints\_video.py}: renderizado de keypoints sobre video para validación.
\end{itemize}

\section{Base de datos}

El sistema utiliza SQLite como motor de persistencia. El esquema almacena
experimentos, usuarios, regiones de interés (ROIs), resultados de clasificación
y rutas a los artefactos generados (video multimodal, trayectorias, timelogs).
Las migraciones se gestionan mediante scripts independientes en \texttt{src/migrations/}
y \texttt{src/db/migrations/}, y se aplican de forma programática durante el arranque
del sistema.

% =============================================================================
\chapter{Pipeline de Procesamiento}
% =============================================================================

\section{Descripción general del flujo}

El pipeline de procesamiento es la secuencia ordenada de transformaciones que
convierte un video de experimento crudo en un conjunto de métricas conductuales
validadas. El orquestador central (\texttt{run\_behavior\_pipeline.py}) gestiona
la ejecución de cada etapa, verifica la frescura de los artefactos intermedios
para evitar recómputos innecesarios y transmite el estado de avance al proceso
padre mediante marcadores estandarizados en la salida estándar.

\section{Etapas del pipeline}

\subsection{Etapa 1 -- Ingesta y validación}

El usuario carga el video a través de la página \texttt{01\_Ingesta\_de\_Video.py}.
El sistema valida el formato, la resolución y la duración del archivo. Los metadatos
del experimento (fecha, identificador del animal, observaciones) quedan registrados
en la base de datos antes de continuar.

\subsection{Etapa 2 -- Extracción de pose}

La extracción de keypoints se ejecuta desde la página \texttt{02\_Keypoints.py},
que lanza \texttt{run\_superanimal.py} en un proceso hijo con la siguiente configuración:

\begin{itemize}[leftmargin=*]
    \item Modelo: DeepLabCut SuperAnimal (\texttt{superanimal\_topviewmouse}).
    \item Entorno Python: \texttt{venv\_310} (Python 3.10, requerimiento de DLC).
    \item Parámetros configurables: tamaño de batch (4--32 frames), rango de tiempo a analizar,
          uso opcional de CPU en caso de indisponibilidad de GPU.
    \item Salida: archivo \texttt{.h5} con coordenadas brutas de keypoints por frame.
\end{itemize}

El proceso emite logs línea a línea y marcadores \texttt{[OUTPUT] clave=valor}
que la interfaz captura para actualizar el estado sin necesidad de polling activo.

\subsection{Etapa 3 -- Filtrado mediante bounding box YOLO}

Una vez disponible el archivo de pose crudo, \texttt{apply\_dlc\_bbox\_constraint.py}
aplica un detector YOLO para localizar el roedor en cada frame. Los keypoints
que caen fuera del bounding box predicho son descartados o corregidos. Esta etapa
reduce artefactos de estimación producidos por oclusiones parciales, sombras o
reflejos en el laberinto.

Los artefactos generados en esta etapa son:
\begin{itemize}[leftmargin=*]
    \item \texttt{*\_bbox\_constrained.h5}: pose filtrada en formato HDF5.
    \item \texttt{*\_bbox\_constrained.csv}: misma información en formato tabular compatible con SimBA.
    \item \texttt{*\_bbox\_constraint.mp4}: video de validación con keypoints y bounding box superpuestos.
\end{itemize}

\subsection{Etapa 4 -- Cómputo de características SimBA}

\texttt{compute\_simba\_features.py} toma el CSV filtrado y las zonas del EPM
configuradas por el usuario, y produce un archivo de características en el
formato que espera el motor de clasificación de SimBA. Las características
incluyen, entre otras: distancias entre keypoints, velocidades instantáneas,
tiempo acumulado por zona, y ángulos de orientación corporal.

\subsection{Etapa 5 -- Inferencia conductual}

Con las características extraídas, el sistema aplica los clasificadores entrenados
para dos conductas objetivo:

\begin{itemize}[leftmargin=*]
    \item \textbf{Thigmotaxis}: tendencia del animal a mantenerse próximo a las paredes
          del laberinto. Se asocia con estados de ansiedad elevada.
    \item \textbf{Grooming}: comportamiento de autolimpieza caracterizado por movimientos
          repetitivos y posturas corporales específicas.
\end{itemize}

Cada clasificador es un modelo de aprendizaje automático (formato \texttt{.sav})
entrenado en SimBA y consumido mediante la API de scikit-learn. Los modelos
residen en \texttt{data/simba\_projects/.../models/generated\_models/}.

\subsection{Etapa 6 -- Generación del video multimodal}

\texttt{generar\_video\_prediccion.py} produce el artefacto visual principal
del sistema: un video donde se superponen sobre los frames originales las
trayectorias del animal, las predicciones conductuales por frame, los
indicadores de zona activa y el marcador temporal. Este video constituye
el material de revisión cualitativa para el investigador.

Los artefactos de esta etapa incluyen:
\begin{itemize}[leftmargin=*]
    \item \texttt{*\_STREAMLIT\_MULTIMODAL.mp4}: video final con anotaciones.
    \item \texttt{*\_trajectory.csv}: coordenadas de trayectoria por frame.
    \item \texttt{*\_GROOMING\_TIMELOG.csv}: intervalos de tiempo con grooming detectado.
    \item \texttt{*\_THIGMOTAXIS\_TIMELOG.csv}: intervalos de tiempo con thigmotaxis detectada.
\end{itemize}

\subsection{Etapa 7 -- Persistencia y visualización}

Los resultados y rutas a artefactos se persisten en la base de datos SQLite.
La página \texttt{05\_Resultados\_y\_Estadisticas.py} los recupera y genera
un dashboard con métricas de tiempo por zona, porcentajes conductuales,
heatmap de posición y reproductor del video multimodal.

% =============================================================================
\chapter{Herramientas y Tecnologías}
% =============================================================================

\section{Entorno de hardware}

El sistema fue desarrollado y validado sobre una estación de trabajo con las
siguientes características de referencia:

\begin{itemize}[leftmargin=*]
    \item Sistema operativo: Windows 11 Pro.
    \item GPU: NVIDIA RTX 5070 Ti Laptop (usada durante etapas DLC e inferencia YOLO).
    \item Almacenamiento: unidad SSD para datos de video y modelos.
\end{itemize}

El sistema puede operar en modo CPU (\texttt{--force-cpu}) con degradación de
velocidad, lo que lo hace funcional en equipos sin acelerador gráfico dedicado,
aunque con tiempos de procesamiento significativamente mayores.

\section{Entornos de ejecución Python}

Debido a las incompatibilidades entre las dependencias de DeepLabCut (Python 3.10)
y las de YOLO/Ultralytics (Python 3.11), el proyecto mantiene dos entornos
virtuales independientes:

\begin{table}[H]
    \centering
    \begin{tabular}{@{}llp{6cm}@{}}
        \toprule
        \textbf{Entorno} & \textbf{Versión} & \textbf{Uso principal} \\
        \midrule
        \texttt{venv\_310} & Python 3.10 & DeepLabCut SuperAnimal, SimBA (features) \\
        \texttt{venv\_311} & Python 3.11 & Ultralytics YOLO, procesamiento de video, interfaz Streamlit \\
        \bottomrule
    \end{tabular}
    \caption{Entornos virtuales del proyecto y sus responsabilidades.}
    \label{tab:envs}
\end{table}

\section{Frameworks y librerías principales}

\begin{table}[H]
    \centering
    \begin{tabular}{@{}lll@{}}
        \toprule
        \textbf{Librería / Framework} & \textbf{Rol en el sistema} & \textbf{Versión objetivo} \\
        \midrule
        Streamlit            & Interfaz web multipágina            & 1.x \\
        DeepLabCut           & Estimación de pose (backend actual) & 3.x \\
        Ultralytics YOLO     & Detección, tracking y pose (futuro) & 8.x / YOLO26 \\
        SimBA                & Features y clasificación conductual & Estable \\
        OpenCV               & Procesamiento y renderizado de video & 4.x \\
        pandas               & Manipulación de datos tabulares     & 2.x \\
        scikit-learn         & Inferencia con modelos \texttt{.sav} & 1.x \\
        SQLite / sqlite3     & Persistencia relacional             & -- \\
        FFmpeg               & Recodificación y recorte de video   & 6.x \\
        python-dotenv        & Gestión de configuración por entorno & -- \\
        \bottomrule
    \end{tabular}
    \caption{Librerías y frameworks principales del proyecto.}
    \label{tab:librerias}
\end{table}

\section{Herramientas de desarrollo}

\begin{itemize}[leftmargin=*]
    \item \textbf{Git}: control de versiones del código fuente.
    \item \textbf{pytest}: suite de pruebas unitarias e de integración.
    \item \textbf{Visual Studio Code}: entorno de desarrollo principal.
    \item \textbf{Overleaf / LaTeX}: documentación técnica formal.
\end{itemize}

% =============================================================================
\chapter{Módulos del Sistema}
% =============================================================================

\section{Módulo de autenticación}

Implementado en \texttt{src/auth.py} y \texttt{src/email\_utils.py}, este módulo
gestiona el ciclo de vida de sesión de los usuarios del sistema. Soporta
verificación en dos factores mediante OTP (One-Time Password) enviado por
correo electrónico. Incluye registro de auditoría de accesos en
\texttt{src/security\_logger.py}.

\section{Módulo de configuración centralizada}

\texttt{src/config.py} define todas las rutas del proyecto de forma relativa
al directorio raíz detectado automáticamente mediante \texttt{Path(\_\_file\_\_).parent.parent}.
Esto elimina rutas absolutas del código fuente y hace el proyecto portable entre
estaciones de trabajo. Incluye detección automática de FFmpeg en el sistema,
variables de entorno opcionales mediante \texttt{.env}, y validación de rutas
críticas al arranque.

\section{Módulo de ingesta de video}

La página \texttt{01\_Ingesta\_de\_Video.py} permite al usuario seleccionar un
video del sistema de archivos local, validar sus metadatos (resolución, duración,
codec) y registrar el experimento en la base de datos antes de continuar al
siguiente paso. El componente de banner de contexto (\texttt{src/video\_context\_banner.py})
mantiene visible la información del video activo en todas las páginas posteriores.

\section{Módulo de extracción de keypoints}

La página \texttt{02\_Keypoints.py} actúa como gestor de la etapa de pose estimation.
Sus responsabilidades son:

\begin{itemize}[leftmargin=*]
    \item Lanzar el proceso de extracción en segundo plano mediante \texttt{run\_with\_live\_log.py}.
    \item Transmitir los logs del subproceso a la interfaz en tiempo real.
    \item Detectar el archivo de pose generado y registrar su ruta en la sesión.
    \item Recuperar el estado de una extracción previa si el usuario regresa a esta página.
    \item Ofrecer la opción experimental de YOLO Pose como backend alternativo.
\end{itemize}

\section{Módulo de configuración de zonas}

\texttt{03\_Configuracion\_Zonas.py} permite al usuario definir las seis
regiones canónicas del EPM (dos brazos abiertos, dos brazos cerrados, área
central y zona perimetral). La interfaz permite ajustar las coordenadas de
cada zona mediante controles deslizantes. El módulo \texttt{src/simba\_roi\_bridge.py}
traduce estas definiciones al formato XML que espera SimBA para sus cálculos
de tiempo por zona.

\section{Módulo de análisis final}

La página \texttt{04\_Analisis\_Final.py} orquesta la ejecución completa del
pipeline desde la interfaz. Invoca \texttt{run\_behavior\_pipeline.py} como
subproceso, recopila los marcadores de salida y persiste los resultados en la
base de datos al finalizar. Maneja también los casos de reanudación parcial,
donde artefactos intermedios ya existentes son reutilizados sin recómputo.

\section{Módulo de resultados}

La página \texttt{05\_Resultados\_y\_Estadisticas.py} construye el dashboard
de visualización final. Carga los datos del experimento desde la base de datos
y genera:

\begin{itemize}[leftmargin=*]
    \item Métricas de tiempo por zona del EPM (porcentajes y segundos absolutos).
    \item Porcentaje de frames clasificados como thigmotaxis y grooming.
    \item Mapa de calor de posición acumulada del animal.
    \item Reproductor del video multimodal generado.
    \item Tablas de timelogs conductuales exportables.
\end{itemize}

\section{Módulo de administración}

\texttt{99\_Admin\_Panel.py} provee funciones de gestión del sistema para
usuarios con rol de administrador: alta y baja de usuarios, revisión del
log de auditoría de seguridad y herramientas de diagnóstico del entorno.

% =============================================================================
\chapter{Funcionamiento del Sistema}
% =============================================================================

\section{Flujo de trabajo desde la perspectiva del usuario}

El sistema está diseñado para guiar al usuario a través de un flujo lineal
de páginas, donde cada etapa depende de la correcta finalización de la anterior.
A continuación se describe el procedimiento completo.

\subsection{Paso 1 -- Acceso al sistema}

El usuario accede a la aplicación Streamlit a través de un navegador web en
la dirección de red local donde está desplegado el servidor. Introduce sus
credenciales en la página de inicio de sesión. Si la autenticación de dos
factores está habilitada, recibirá un código OTP en su correo institucional
que deberá ingresar para completar el acceso.

\subsection{Paso 2 -- Ingesta del video}

En la página de ingesta, el usuario:
\begin{enumerate}[label=\alph*.]
    \item Selecciona el archivo de video desde su equipo local.
    \item Revisa los metadatos detectados automáticamente (duración, resolución, fps).
    \item Registra los datos del experimento (identificador del animal, fecha, observaciones).
    \item Confirma la ingesta para continuar.
\end{enumerate}

\subsection{Paso 3 -- Extracción de keypoints}

El usuario inicia el proceso de extracción desde la página de Keypoints.
El sistema lanza DeepLabCut en segundo plano y muestra los logs de avance en
tiempo real. Dependiendo de la duración del video y el hardware disponible,
este proceso puede durar entre 30 minutos y varias horas. Una vez finalizado,
el sistema notifica al usuario y habilita la continuación al siguiente paso.

\subsection{Paso 4 -- Configuración de zonas del EPM}

El usuario ajusta las coordenadas de las zonas del EPM para el video actual.
El sistema ofrece una configuración predeterminada basada en plantillas, que
el usuario puede modificar mediante controles en la interfaz. Las zonas se
guardan y sincronizan automáticamente con el proyecto SimBA activo.

\subsection{Paso 5 -- Análisis final}

El usuario lanza el análisis final desde la página correspondiente. El sistema
ejecuta las etapas de filtrado YOLO, cómputo de características SimBA e
inferencia conductual de forma automática y secuencial. Al finalizar, el
video multimodal y las métricas quedan disponibles en la base de datos.

\subsection{Paso 6 -- Consulta de resultados}

El usuario accede al dashboard de resultados para revisar las métricas
cuantitativas del experimento, visualizar el video multimodal y descargar
los archivos de timelogs conductuales para análisis estadístico posterior.

% =============================================================================
\chapter{Validación y Pruebas}
% =============================================================================

\section{Estrategia de validación}

La validación del sistema se aborda en dos dimensiones complementarias: la
validación técnica del software (corrección funcional de los módulos) y la
validación científica de los modelos (precisión de la clasificación conductual
respecto a una referencia experta).

\section{Pruebas del software}

La suite de pruebas del proyecto reside en el directorio \texttt{tests/} e
incluye siete archivos con un total de 387 líneas de código de prueba. La
cobertura actual se concentra en:

\begin{itemize}[leftmargin=*]
    \item Autenticación de usuarios y gestión de sesiones.
    \item Operaciones de lectura y escritura en la base de datos.
    \item Lógica de análisis conductual base (\texttt{analysis\_logic.py}).
    \item Validación de rutas y configuración centralizada.
\end{itemize}

Las pruebas se ejecutan con pytest y pueden lanzarse desde la raíz del proyecto.
La cobertura del pipeline de procesamiento (etapas DLC, filtrado YOLO, bridge
SimBA) está identificada como área de mejora prioritaria para la segunda mitad
de TT2.

\section{Validación del pipeline con datos reales}

El pipeline completo ha sido ejecutado satisfactoriamente con videos reales
del experimento EPM. Una corrida documentada procesó el video
\texttt{R6YB15\_01mar24.mp4} (311 segundos de experimento) en un tiempo
efectivo de 3 horas 19 minutos utilizando DeepLabCut con GPU RTX 5070 Ti.

Este tiempo constituye el indicador de línea base más importante del proyecto
y justifica directamente la investigación de alternativas de pose más eficientes
como objetivo central de TT2.

\section{Validación de clasificadores conductuales}

Los modelos de thigmotaxis y grooming fueron entrenados en SimBA con anotaciones
manuales de video de referencia. La evaluación de estos modelos se realiza
comparando las predicciones automáticas del sistema con el criterio de un
observador experto en comportamiento animal. Esta validación científica se
encuentra en progreso y constituye uno de los frentes de trabajo activos de
la segunda etapa.

% =============================================================================
\chapter{Limitaciones del Sistema}
% =============================================================================

\section{Limitaciones técnicas}

\subsection{Latencia del backend de pose}

La limitación más significativa del sistema en su estado actual es el tiempo
de procesamiento requerido por DeepLabCut SuperAnimal para la extracción de
keypoints. Los tiempos medidos oscilan entre 3 y 4.5 horas por video de 5
minutos en hardware de referencia. Este cuello de botella impide la iteración
experimental ágil y dificulta el análisis de lotes de videos.

\subsection{Dependencia de entornos duales}

La necesidad de mantener dos entornos Python independientes (\texttt{venv\_310}
para DLC y \texttt{venv\_311} para YOLO/Streamlit) añade complejidad operativa
al sistema. La comunicación entre etapas se realiza mediante subprocesos con
marcadores de texto en la salida estándar, lo que, si bien funciona correctamente,
resulta frágil ante cambios de formato de log en versiones futuras de DLC.

\subsection{Esquema de base de datos con migraciones en runtime}

Algunas columnas del esquema (como \texttt{trajectory\_path}) se agregan
mediante sentencias \texttt{ALTER TABLE ... ADD COLUMN IF NOT EXISTS} ejecutadas
en tiempo de ejecución de la aplicación, en lugar de ser gestionadas por un
sistema formal de migraciones. Esto puede producir comportamientos inesperados
en entornos de producción multi-instancia.

\subsection{Escenario experimental único}

El dataset de entrenamiento de los modelos conductuales fue construido con
videos de un único escenario experimental (mismo aparato EPM, misma cámara,
condiciones de iluminación similares). Esto puede limitar la generalización
de los clasificadores ante variaciones en el equipamiento o las condiciones
de captura.

\subsection{Cobertura de pruebas del pipeline}

Las pruebas automatizadas no cubren aún las etapas más críticas del pipeline
(extracción de pose, filtrado YOLO, cómputo de características). Esto reduce
la capacidad de detectar regresiones al introducir cambios en estos módulos.

\section{Limitaciones científicas}

La clasificación de grooming depende de micro-posturas corporales que pueden
no estar suficientemente representadas con el esquema actual de 8 keypoints.
La validez de los clasificadores entrenados está condicionada a la calidad y
diversidad de las anotaciones de referencia, cuya construcción formal está
en curso.

% =============================================================================
\chapter{Trabajo a Futuro}
% =============================================================================

\section{Migración a YOLO26 Pose}

La tarea de mayor prioridad para TT2 es la integración de YOLO26 Pose como
backend alternativo de extracción de pose. La motivación central es reducir
el tiempo de inferencia de varias horas a un rango de minutos por video, lo
que habilitaría un flujo de trabajo experimental sustancialmente más ágil.

La ruta de integración planificada es:

\begin{enumerate}[label=\arabic*.]
    \item Completar y validar el dataset de entrenamiento de pose:
          857 imágenes anotadas, 5 videos fuente, 1 escenario, 8 keypoints por pose.
    \item Entrenar el modelo YOLO26 Pose y evaluar métricas de precisión de keypoints.
    \item Desarrollar el adaptador de salida que convierta las predicciones YOLO Pose
          al formato CSV de entrada esperado por \texttt{compute\_simba\_features.py}.
    \item Integrar el nuevo backend en \texttt{pages/02\_Keypoints.py} como opción
          seleccionable al lado del backend DLC actual.
    \item Ejecutar un benchmark comparativo formal entre DLC y YOLO26 Pose en términos
          de velocidad, estabilidad de keypoints e impacto en la clasificación conductual.
    \item Conservar YOLO11 Pose como baseline de comparación de velocidad y estabilidad.
\end{enumerate}

\section{Refactorización de la arquitectura de backends de pose}

Para soportar múltiples backends de forma mantenible, se propone diseñar una
interfaz formal (\textit{abstract base class}) que defina el contrato que
cualquier backend de pose debe cumplir para integrarse con el resto del sistema.
Esto eliminaría la necesidad de bloques condicionales dispersos en el pipeline
y facilitaría la incorporación futura de nuevos backends.

\section{Mejoras en la cobertura de pruebas}

\begin{itemize}[leftmargin=*]
    \item Pruebas de integración del pipeline con un clip corto de referencia
          y salidas esperadas predefinidas (smoke test).
    \item Pruebas de las etapas de filtrado YOLO, cómputo de características
          y bridge de zonas.
    \item Pruebas de regresión que detecten automáticamente cambios en la
          distribución de predicciones conductuales.
\end{itemize}

\section{Gestión formal de migraciones de base de datos}

Migrar las modificaciones de esquema en runtime a un sistema de migraciones
versionado (por ejemplo, mediante Alembic o un mecanismo propio basado en
un registro de versiones aplicadas), para garantizar que el esquema sea
reproducible y auditaable entre instalaciones.

\section{Ampliación del dataset conductual}

\begin{itemize}[leftmargin=*]
    \item Incorporar videos de experimentos con condiciones de iluminación
          y ángulos de cámara variados para mejorar la generalización.
    \item Reforzar la representación de secuencias de grooming en el conjunto
          de entrenamiento, dado el interés científico específico en esta conducta
          y la dependencia de sus micro-posturas en una representación de keypoints
          de alta calidad.
    \item Establecer una verdad terreno formal para pose y comportamiento
          mediante anotación experta de un subconjunto de video de validación.
\end{itemize}

\section{Actualización de la documentación}

Actualizar el contenido de \texttt{Home.py}, el diagrama \texttt{pipeline\_TT2026.html}
y los manuales de pruebas para reflejar el estado real del sistema, eliminando
referencias a módulos no conectados actualmente al flujo principal.

% =============================================================================
\chapter{Anexos Técnicos}
% =============================================================================

\section{Estructura del repositorio}

El repositorio del proyecto tiene la siguiente organización de alto nivel:

\begin{lstlisting}[language=bash, caption={Estructura de directorios del repositorio.}]
TT_Ratones_2026/
├── Home.py                    # Página principal de la app Streamlit
├── pages/                     # Páginas de la aplicación (9 archivos, ~5230 líneas)
├── src/                       # Lógica central del sistema (~82 archivos, ~10752 líneas)
│   ├── config.py              # Configuración centralizada
│   ├── auth.py                # Autenticación y OTP
│   ├── simba_roi_bridge.py    # Sincronización de zonas con SimBA
│   ├── db/                    # Capa de acceso a base de datos
│   ├── migrations/            # Scripts de migración de esquema
│   └── scripts/               # Scripts del pipeline y herramientas
├── tests/                     # Suite de pruebas (7 archivos, ~387 líneas)
├── data/                      # Modelos, proyectos SimBA y artefactos
├── videos_data/               # Videos de experimento
├── logs/                      # Logs del sistema y del pipeline
├── reportes/                  # Documentación técnica del proyecto
├── assets/                    # Recursos estáticos (logos, imágenes)
├── launcher.bat               # Lanzador del sistema en Windows
└── start_services.py          # Script de arranque de servicios
\end{lstlisting}

\section{Scripts de utilidad relevantes}

Además de los scripts del pipeline principal, el directorio \texttt{src/scripts/}
contiene una colección de utilidades de apoyo al desarrollo y validación:

\begin{table}[H]
    \centering
    \begin{tabular}{@{}lp{8.5cm}@{}}
        \toprule
        \textbf{Script} & \textbf{Propósito} \\
        \midrule
        \texttt{entrenar\_yolo\_local.py}      & Entrena un detector YOLO con el dataset local \\
        \texttt{solo\_yolo\_tracker.py}         & Tracking puro con YOLO sin estimación de pose \\
        \texttt{verify\_gpu.py}                & Verifica la disponibilidad y estado de la GPU \\
        \texttt{benchmark\_gpu.py}             & Benchmark de rendimiento del acelerador \\
        \texttt{trim\_test\_video.py}           & Recorta clips de prueba desde videos de experimento \\
        \texttt{render\_dlc\_keypoints\_video.py} & Renderiza keypoints DLC sobre video para revisión \\
        \texttt{prune\_simba\_project.py}       & Limpia artefactos intermedios del proyecto SimBA \\
        \texttt{run\_masivo\_dlc.py}            & Procesamiento en lote de múltiples videos con DLC \\
        \bottomrule
    \end{tabular}
    \caption{Scripts de utilidad del directorio \texttt{src/scripts/}.}
    \label{tab:scripts_util}
\end{table}

\section{Configuración del entorno}

El sistema requiere las siguientes variables de entorno, definidas en un
archivo \texttt{.env} en la raíz del proyecto:

\begin{lstlisting}[caption={Variables de entorno del sistema.}]
# Modelo YOLO activo (ruta relativa desde la raíz del proyecto)
YOLO_MODEL=yolo_tracker.pt

# Configuración SMTP para envío de OTP
SMTP_HOST=smtp.example.com
SMTP_PORT=587
SMTP_USER=usuario@example.com
SMTP_PASSWORD=contraseña

# Ruta opcional a FFmpeg (si no está en PATH)
FFMPEG_PATH=C:/ffmpeg/bin/ffmpeg.exe
\end{lstlisting}

\section{Instrucciones de arranque}

El sistema puede iniciarse mediante el archivo \texttt{launcher.bat} en Windows,
o directamente con:

\begin{lstlisting}[language=bash, caption={Arranque manual del sistema.}]
# Activar el entorno virtual principal
venv_311\Scripts\activate

# Lanzar la aplicación Streamlit
streamlit run Home.py
\end{lstlisting}

\section{Métricas de tamaño del código fuente}

Al inicio de la fase TT2, el repositorio activo contiene:

\begin{table}[H]
    \centering
    \begin{tabular}{@{}lrr@{}}
        \toprule
        \textbf{Directorio} & \textbf{Archivos} & \textbf{Líneas de código} \\
        \midrule
        \texttt{pages/}  & 9  & 5,230  \\
        \texttt{src/}    & 82 & 10,752 \\
        \texttt{tests/}  & 7  & 387    \\
        \midrule
        \textbf{Total}   & \textbf{98} & \textbf{16,369} \\
        \bottomrule
    \end{tabular}
    \caption{Métricas de tamaño del código fuente activo al inicio de TT2.}
    \label{tab:loc}
\end{table}

\end{document}
